% Main document: graphicx and tikz; serif typography follows the main document.
% The same augmentations are applied to the real image and prototype thumbnails.
\begin{figure*}[t]
\centering
\begingroup
\def\MIPAImageDirectory{images}
\def\MIPACenterImage{\MIPAImageDirectory/real.png}
% Clockwise from top: filename / first-view probability / second-view probability.
% Illustrative probabilities; each column sums to 100.
\def\MIPAPrototypeData{%
  prototype-01/3/5,
  prototype-02/2/1,
  prototype-03/27/34,
  prototype-04/3/4,
  prototype-05/38/28,
  prototype-06/6/8,
  prototype-07/4/3,
  prototype-08/2/1,
  prototype-09/3/5,
  prototype-10/12/11%
}
\def\MIPAMaxPercent{40}
\def\MIPAZoom{1.3}
\definecolor{MIPAReal}{RGB}{91,132,180}
\definecolor{MIPASynthetic}{RGB}{199,106,100}
\definecolor{MIPABlueFill}{RGB}{217,230,244}
\definecolor{MIPARedFill}{RGB}{248,222,218}
\definecolor{MIPAText}{RGB}{40,44,49}
\definecolor{MIPAMuted}{RGB}{91,101,111}
\colorlet{MIPAGrid}{MIPAMuted!17!white}
\colorlet{MIPARim}{MIPAMuted!30!white}
\colorlet{MIPAFrame}{MIPASynthetic!60!white}
% Arguments: PNG path, center x, center y, square side in bp, view number.
% Rectangular images are center-cropped to square with their aspect ratio kept.
\newsavebox{\MIPAImageBox}
\newcommand{\MIPAImage}[5]{%
  \begin{scope}[shift={(#2,#3)}]
    \pgfmathsetmacro{\MIPAHalf}{#4/2}
    \clip (-\MIPAHalf,-\MIPAHalf) rectangle (\MIPAHalf,\MIPAHalf);
    \pgfmathsetlengthmacro{\MIPASide}{#4*1bp}
    \ifnum#5=2\relax
      \pgfmathsetlengthmacro{\MIPASide}{#4*\MIPAZoom*1bp}
    \fi
    \sbox{\MIPAImageBox}{\includegraphics{#1}}%
    \ifdim\wd\MIPAImageBox>\ht\MIPAImageBox
      \sbox{\MIPAImageBox}{\includegraphics[height=\MIPASide]{#1}}%
    \else
      \sbox{\MIPAImageBox}{\includegraphics[width=\MIPASide]{#1}}%
    \fi
    \ifnum#5=1\relax
      \node[rotate=-90,transform shape,inner sep=0pt,outer sep=0pt]
        at (0,0) {\usebox{\MIPAImageBox}};
    \else
      \node[xscale=-1,transform shape,inner sep=0pt,outer sep=0pt]
        at (0,0) {\usebox{\MIPAImageBox}};
    \fi
  \end{scope}%
}

\resizebox{\linewidth}{!}{%
\begin{tikzpicture}[x=1bp,y=1bp,inner sep=0pt,outer sep=0pt,
  every node/.style={font=\rmfamily}]
  \path[use as bounding box] (0,-268) rectangle (504,0);
  \fill[white] (0,-268) rectangle (504,0);
  \foreach \MIPAView/\MIPACenterX in {1/126,2/378} {
    \ifnum\MIPAView=1
      \colorlet{MIPAAccent}{MIPAReal}
      \colorlet{MIPABarFill}{MIPABlueFill}
      \def\MIPAPanel{a}\def\MIPATitle{Rotated view}
      \def\MIPASubtitle{$90^{\circ}$ clockwise}
    \else
      \colorlet{MIPAAccent}{MIPASynthetic}
      \colorlet{MIPABarFill}{MIPARedFill}
      \def\MIPAPanel{b}\def\MIPATitle{Flipped and cropped view}
      \def\MIPASubtitle{Horizontal flip, $1.3\!\times$ zoom}
    \fi
    \node[anchor=base,text=MIPAText,font=\rmfamily\bfseries\fontsize{10.5bp}{12bp}\selectfont]
      at (\MIPACenterX,-12) {(\MIPAPanel)\enspace\MIPATitle};
    \node[anchor=base,text=MIPAMuted,font=\rmfamily\fontsize{8.5bp}{10bp}\selectfont]
      at (\MIPACenterX,-26) {\MIPASubtitle};
    \begin{scope}[shift={(\MIPACenterX,-143)}]
      % Keep the common radial scale, but quiet the scaffolding.
      \draw[MIPARim,line width=.45bp] (0,0) circle[radius=90bp];
      \foreach \MIPAFraction in {.5,1} {
        \pgfmathsetmacro{\MIPARadius}{35.5+29.5*\MIPAFraction}
        \draw[MIPAGrid,line width=.4bp,dash pattern=on 1bp off 2bp]
          (0,0) circle[radius=\MIPARadius bp];
      }
      \foreach \MIPAFile/\MIPAFirst/\MIPASecond
        [count=\MIPAIndex from 0] in \MIPAPrototypeData {
        \ifnum\MIPAView=1\let\MIPAProbability\MIPAFirst
        \else\let\MIPAProbability\MIPASecond\fi
        \pgfmathsetmacro{\MIPAAngle}{90-36*\MIPAIndex}
        \pgfmathsetmacro{\MIPAStart}{\MIPAAngle-7}
        \pgfmathsetmacro{\MIPAEnd}{\MIPAAngle+7}
        \pgfmathsetmacro{\MIPAOuter}{35.5+29.5*\MIPAProbability/\MIPAMaxPercent}
        \colorlet{MIPABar}{MIPABarFill}
        \colorlet{MIPABarEdge}{MIPAAccent!80!white}
        \draw[MIPAGrid,line width=.35bp] (\MIPAAngle:35.5) -- (\MIPAAngle:90);
        \path[fill=MIPABar,draw=MIPABarEdge,line width=.35bp]
          (\MIPAStart:35.5) -- (\MIPAStart:\MIPAOuter)
          arc[start angle=\MIPAStart,end angle=\MIPAEnd,radius=\MIPAOuter bp]
          -- (\MIPAEnd:35.5)
          arc[start angle=\MIPAEnd,end angle=\MIPAStart,radius=35.5bp] -- cycle;
      }
      \filldraw[fill=white,draw=MIPARim,line width=.6bp] (0,0) circle[radius=35.5bp];
      \filldraw[fill=white,draw=MIPAReal,line width=.8bp]
        (-24.6,-24.6) rectangle (24.6,24.6);
      \MIPAImage{\MIPACenterImage}{0}{0}{48}{\MIPAView}
      \foreach \MIPAFile/\MIPAFirst/\MIPASecond
        [count=\MIPAIndex from 0] in \MIPAPrototypeData {
        \ifnum\MIPAView=1\let\MIPAProbability\MIPAFirst
        \else\let\MIPAProbability\MIPASecond\fi
        \pgfmathsetmacro{\MIPAAngle}{90-36*\MIPAIndex}
        \colorlet{MIPABorder}{MIPAFrame}\colorlet{MIPALabel}{MIPAMuted}
        \def\MIPABorderWidth{.5bp}\let\MIPAWeight\mdseries
        \ifdim\MIPAProbability pt>20pt
          \colorlet{MIPABorder}{MIPAAccent}\colorlet{MIPALabel}{MIPAText}
          \def\MIPABorderWidth{.85bp}\let\MIPAWeight\bfseries
        \fi
        \begin{scope}[shift={(\MIPAAngle:90)}]
          \fill[white] (-16,-16) rectangle (16,16);
          \filldraw[fill=white,draw=MIPABorder,line width=\MIPABorderWidth]
            (-14.7,-14.7) rectangle (14.7,14.7);
          \MIPAImage{\MIPAImageDirectory/\MIPAFile.png}{0}{0}{28}{\MIPAView}
          \node[anchor=north,fill=white,inner xsep=2bp,inner ysep=1.3bp,
            text=MIPALabel,font=\rmfamily\MIPAWeight\fontsize{9bp}{10bp}\selectfont]
            at (0,-16.1) {\MIPAProbability\%};
        \end{scope}
      }
    \end{scope}
  }
\end{tikzpicture}%
}
\endgroup
\caption{\textbf{Illustrative soft assignments across views.} Two augmented views
of one real image (center) are compared with the same ten synthetic prototypes
(rim). Light blue and red bars encode assignment probability on a shared 0--40\% scale.
All probabilities are illustrative.}
\label{fig:prototype-assignments}
\end{figure*}
